% !TEX program = lualatex
\documentclass[11pt, a4paper]{article}

\usepackage{fontspec}
\setmainfont{TeX Gyre Pagella}

\usepackage{amsmath, amssymb, amsthm}
\usepackage{geometry}
\usepackage{microtype}
\usepackage{setspace}
\usepackage{titlesec}
\usepackage{hyperref}
\usepackage{mdframed}
\usepackage{babel}

\geometry{
  top=1.25in,
  bottom=1.25in,
  left=1.35in,
  right=1.35in
}

\onehalfspacing

\hypersetup{
  colorlinks=true,
  linkcolor=black,
  citecolor=black,
  urlcolor=black
}

\newtheoremstyle{named}
  {12pt}{12pt}{\itshape}{}{\bfseries}{.}{.5em}{}
\theoremstyle{named}
\newtheorem{proposicao}{Proposição}
\newtheorem{definicao}{Definição}

\newmdenv[
  linewidth=0.8pt,
  innertopmargin=10pt,
  innerbottommargin=10pt,
  innerleftmargin=14pt,
  innerrightmargin=14pt,
  skipabove=14pt,
  skipbelow=14pt
]{caixachave}

\titleformat{\section}{\large\bfseries}{\thesection.}{0.5em}{}
\titleformat{\subsection}{\normalsize\bfseries}{\thesubsection.}{0.5em}{}

\title{%
  \textbf{Computação Além dos Dados}\\[0.4em]
  \large Restrições, Acessibilidade e Memória como\\
  Fundamentos de Sistemas Complexos
}
\author{Flyxion}
\date{}

% ---------------------------------------------------------------
\begin{document}

\maketitle
\thispagestyle{empty}

\begin{abstract}
Grande parte da Ciência da Computação moderna é construída sobre uma
suposição aparentemente óbvia: primeiro existem os dados, depois os
algoritmos operam sobre eles. Este ensaio propõe uma inversão parcial
dessa perspectiva. Em vez de tratar objetos, dados ou estados como
elementos fundamentais, o foco passa para as restrições, trajetórias
possíveis e espaços de acessibilidade que tornam certos estados
alcançáveis e outros impossíveis. Essa mudança de ênfase não é apenas
filosófica. Ela tem consequências diretas para a forma como projetamos
sistemas de IA, modelamos memória computacional, entendemos arquiteturas
distribuídas e concebemos a cognição. O ensaio apresenta cinco
frameworks interconectados — CLIO, RSVP, MEM\textbar{}8, HYDRA e TARTAN
— como instâncias de um programa teórico unificado: a ideia de que
restrição precede conteúdo, que acessibilidade é mais fundamental que
distância, e que a pergunta mais produtiva sobre qualquer sistema
complexo não é \emph{o que existe?} mas \emph{o que ainda é possível?}
\end{abstract}

\newpage
\tableofcontents
\newpage

% ---------------------------------------------------------------
\section{A Suposição Invisível}

Todo estudante de Ciência da Computação aprende, cedo, a pensar em
termos de dados e operações. Existe um array; aplicamos um algoritmo de
ordenação. Existe um grafo; executamos uma busca em largura. Existe um
conjunto de exemplos rotulados; treinamos um classificador. A estrutura
geral é sempre a mesma: o mundo fornece objetos, e nós os processamos.

Essa suposição é tão fundamental que raramente é examinada. Ela aparece
na terminologia básica — variável, valor, estado, dado — e organiza
silenciosamente a forma como pensamos sobre quase tudo na área.

Mas existe uma outra forma de pensar. Em vez de começar pelos objetos e
perguntar o que podemos fazer com eles, podemos começar pelas restrições
e perguntar quais objetos elas tornam possíveis. Em vez de começar pelos
estados e perguntar como transicionamos entre eles, podemos começar
pelos espaços de estados e perguntar quais regiões são acessíveis a
partir de onde estamos.

Essa inversão não é apenas uma curiosidade filosófica. Ela produz
perguntas diferentes, modelos diferentes, e, argumentamos aqui, uma
compreensão mais profunda de vários dos problemas mais difíceis da
computação moderna: o problema da generalização em aprendizado de
máquina, o problema da consistência em sistemas distribuídos, o problema
da persistência em memória computacional, o problema da emergência em
sistemas cognitivos.

O objetivo deste ensaio é apresentar um conjunto de frameworks
teóricos que partem dessa inversão e exploram suas consequências em
diferentes domínios. Os frameworks são independentes o suficiente para
serem estudados separadamente, mas compartilham um núcleo conceitual
comum que os torna partes de um programa teórico unificado.

% ---------------------------------------------------------------
\section{O Princípio Central: Restrição Antes do Conteúdo}

\subsection{O Jogo de Xadrez}

Considere um jogo de xadrez. As peças importam, mas as regras que
definem os movimentos possíveis importam ainda mais. Sem regras, não
existe jogo — existe apenas um conjunto de objetos de madeira ou
plástico sobre uma superfície quadriculada. O jogo emerge das restrições,
não dos componentes.

Isso não é trivial. Significa que a identidade de uma peça — o que ela
\emph{é}, funcionalmente — é completamente determinada pelas restrições
que governam seu comportamento. Um rei que pudesse se mover como uma
rainha não seria mais um rei. Trocar a regra é trocar o objeto, mesmo
que a peça física permaneça a mesma.

A mesma lógica aparece em toda a computação:

Em compiladores, a gramática formal define o espaço de programas
sintaticamente válidos. Um programa não é apenas uma sequência de
caracteres; é uma sequência que satisfaz as restrições impostas pela
gramática. A gramática precede e constitui o programa.

Em bancos de dados relacionais, o esquema define quais registros podem
existir. Um banco de dados não é apenas uma coleção de dados; é uma
coleção de dados que satisfaz um conjunto de restrições de integridade.
Violar essas restrições não é apenas um erro técnico; é uma contradição
com a própria definição do que o banco de dados representa.

Em redes neurais, a arquitetura define o espaço de funções que a rede
pode aprender. Antes de ver qualquer exemplo, a arquitetura já determina
quais representações são possíveis e quais são impossíveis. A restrição
arquitetural precede os dados de treinamento.

Em sistemas operacionais, o modelo de permissões define quais ações
cada processo pode executar. O conceito de segurança não é um add-on
adicionado após o sistema ser construído; é a estrutura de restrições
sobre a qual o sistema opera.

\subsection{A Proposição Central}

\begin{caixachave}
\textbf{Princípio de Restrição.} As restrições que definem um sistema
não são detalhes de implementação adicionados para evitar erros. São a
principal fonte de organização do sistema — a estrutura que determina
quais estados existem, quais transições são possíveis, e quais objetivos
podem ser perseguidos.
\end{caixachave}

Essa perspectiva tem uma consequência imediata: para entender um sistema
complexo, a primeira pergunta não é \emph{quais elementos ele contém?}
mas \emph{quais restrições o organizam?}

\subsection{Restrições como Informação}

Há uma conexão profunda entre restrições e informação que vale a pena
tornar explícita. Shannon definiu informação como a quantidade que reduz
incerteza. Uma mensagem é informativa na medida em que elimina
alternativas possíveis.

Restrições fazem exatamente isso. Uma restrição elimina estados
possíveis. Quanto mais restritivo é um sistema, menor é o espaço de
estados acessíveis — e, portanto, mais informação está implicitamente
codificada na estrutura de restrições.

Essa observação tem implicações práticas. Quando projetamos um sistema,
cada decisão de design que reduz o espaço de estados possíveis é também
uma decisão que codifica informação sobre o propósito do sistema.
Um sistema sem restrições é um sistema sem propósito — e um sistema
sem propósito é computacionalmente inexpressivo.

\begin{proposicao}[Equivalência Restrição-Informação]
Seja $\mathcal{S}$ o espaço completo de estados de um sistema e
$\mathcal{C}$ o subconjunto de estados admissíveis sob um conjunto
de restrições $R$. A quantidade de informação codificada por $R$ é
proporcional a $\log |\mathcal{S}| - \log |\mathcal{C}|$: a diferença
entre o espaço irrestrito e o espaço restrito. Restrições mais fortes
codificam mais informação.
\end{proposicao}

% ---------------------------------------------------------------
\section{CLIO: Inteligência como Navegação em Espaços de Projeção}

\subsection{O Problema da Observação Parcial}

Sistemas inteligentes raramente observam a realidade completa. Eles
recebem projeções parciais do mundo: pixels em vez de objetos, tokens
em vez de significados, sinais sensoriais em vez de estados físicos.

Formalmente, seja $X$ o espaço completo de estados do mundo e $M$ o
espaço de observações disponíveis para o sistema. Existe uma função de
projeção:
\[
  \pi: X \rightarrow M
\]
que mapeia estados do mundo em observações. Em geral, $\pi$ não é
injetora: múltiplos estados do mundo podem produzir a mesma observação.

O problema central da inteligência, nessa perspectiva, é a inversão
dessa projeção: dado $m \in M$, qual é o conjunto
\[
  \pi^{-1}(m) = \{x \in X \mid \pi(x) = m\}
\]
de estados do mundo consistentes com a observação?

\subsection{Incerteza como Estrutura}

A pré-imagem $\pi^{-1}(m)$ representa a incerteza estrutural do sistema:
a ambiguidade que é matematicamente inerente à projeção, independente de
qualquer limitação de processamento. Um sistema que observa $m$ não pode
distinguir entre nenhum dos estados em $\pi^{-1}(m)$ apenas com base
nessa observação.

Isso tem várias consequências importantes:

\textbf{Para aprendizado de máquina:} o problema de generalização é
fundamentalmente o problema de inverter uma projeção. Um modelo treinado
em exemplos viu apenas uma amostra do espaço de estados do mundo. A
generalização é a capacidade de reconstruir a estrutura do espaço de
estados a partir dessas amostras.

\textbf{Para sistemas de percepção:} a ambiguidade perceptual não é
um bug que pode ser eliminado com mais dados ou melhor hardware. É uma
propriedade estrutural da relação entre sensores e mundo.

\textbf{Para raciocínio sob incerteza:} a incerteza bayesiana é uma
forma de representar $\pi^{-1}(m)$ através de distribuições de
probabilidade sobre estados do mundo.

\subsection{Inteligência como Redução Eficiente de Ambiguidade}

O framework CLIO define inteligência operacionalmente como a capacidade
de reduzir o espaço $\pi^{-1}(m)$ de maneira eficiente, usando
restrições adicionais derivadas do contexto, da memória e do raciocínio.

\begin{definicao}[Inteligência CLIO]
Um sistema é inteligente na medida em que, dado $m \in M$, consegue
identificar um subconjunto $\mathcal{A}(m) \subseteq \pi^{-1}(m)$
de estados plausíveis tal que $|\mathcal{A}(m)| \ll |\pi^{-1}(m)|$,
usando restrições derivadas de contexto e memória, sem acesso direto
a $X$.
\end{definicao}

Essa definição unifica vários conceitos que normalmente aparecem
separados em IA:

\textbf{Raciocínio} é a aplicação de restrições lógicas para reduzir
$\pi^{-1}(m)$.

\textbf{Aprendizado} é a aquisição de restrições que tornam a redução
mais eficiente em novos contextos.

\textbf{Memória} é o armazenamento de restrições derivadas de
experiências anteriores.

\textbf{Generalização} é a transferência de restrições para contextos
não observados.

\subsection{Falha de Projeção e Paradoxo}

Uma consequência importante do framework CLIO é uma explicação natural
para certos tipos de paradoxo e ilusão cognitiva.

\begin{definicao}[Falha de Projeção]
Uma falha de projeção ocorre quando um sistema trata a projeção $m$
como se fosse o estado completo $x$ — quando confunde o mapa com o
território.
\end{definicao}

Muitos erros sistemáticos de raciocínio humano e de sistemas de IA
podem ser entendidos como falhas de projeção. Vieses cognitivos surgem
quando o sistema usa heurísticas que funcionam bem para a maioria dos
estados em $\pi^{-1}(m)$ mas falham para estados menos típicos.
Alucinações em LLMs surgem quando o modelo trata sua representação
interna como se fosse o estado real do mundo.

Paradoxos semânticos surgem quando uma descrição $m$ é tratada como
referência direta a um estado $x$, quando na verdade $\pi^{-1}(m)$
contém múltiplos estados com propriedades incompatíveis. Ao restaurar
o contexto — ao especificar em qual estado de $\pi^{-1}(m)$ estamos
— o paradoxo frequentemente dissolve.

% ---------------------------------------------------------------
\section{RSVP: Sistemas Complexos como Processos de Relaxamento}

\subsection{Da Cosmologia à Computação}

O framework RSVP foi desenvolvido originalmente como uma teoria
cosmológica alternativa, mas sua estrutura matemática possui uma
interpretação computacional direta e independente da motivação
original.

A ideia central é que sistemas complexos evoluem para estados que
reduzem tensões estruturais internas. Em vez de imaginar um sistema
como um conjunto estático de objetos em estados fixos, RSVP o descreve
como um processo dinâmico de reorganização contínua de restrições.

\subsection{Potencial Institucional e Descida de Gradiente}

Formalmente, seja $x(t)$ o estado de um sistema no tempo $t$ e
$V(x, t)$ uma função potencial que representa tensões estruturais
internas — inconsistências, restrições violadas, gradientes não
resolvidos. O sistema evolui aproximadamente por descida de gradiente:
\[
  \frac{dx}{dt} = -\nabla_x V(x,t) + \xi(t),
\]
onde $\xi(t)$ representa perturbações externas.

Cientistas da computação reconhecerão essa estrutura imediatamente:
é a mesma de descida de gradiente estocástica em aprendizado de máquina,
de otimização variacional em modelos generativos, e de equações de
difusão em sistemas físicos.

A diferença de perspectiva é que RSVP trata a função $V$ não como
uma função objetivo escolhida pelo projetista, mas como uma propriedade
emergente da estrutura de restrições do sistema. A ``otimização'' é
interna e não requer um objetivo externo explícito.

\subsection{Atratores e Configurações Estáveis}

Um atrator é uma região do espaço de estados para a qual o sistema
converge naturalmente. Na perspectiva RSVP, atratores correspondem a
configurações que minimizam tensões estruturais localmente.

\begin{proposicao}[Atratores como Mínimos de Tensão]
As configurações estáveis de um sistema complexo correspondem a mínimos
locais da função de tensão $V(x,t)$. Uma perturbação desestrutura o
sistema ao aumentar $V$ além de um limiar crítico, tornando o mínimo
local instável — quando a Hessiana $\nabla_x^2 V(x^*,t)$ perde
positividade definida.
\end{proposicao}

Esse framework tem aplicações diretas para entender:

\textbf{Treinamento de redes neurais:} o processo de treinamento é a
busca por mínimos de uma função de perda. A paisagem de perda tem
múltiplos mínimos locais, e diferentes configurações de hiperparâmetros
correspondem a diferentes funções de tensão.

\textbf{Sistemas distribuídos:} a consistência eventual em sistemas
distribuídos pode ser entendida como convergência para um atrator
global a partir de estados localmente inconsistentes.

\textbf{Modelos de linguagem:} a geração de texto pode ser vista como
navegação em um espaço de estados linguísticos com uma função de tensão
implícita definida pelos dados de treinamento.

\subsection{Entropia, Informação e Geometria}

Uma das contribuições conceituais mais ricas de RSVP é a
reinterpretação de entropia e informação como propriedades geométricas
da dinâmica do sistema.

Em vez de tratar entropia como uma quantidade escalar associada a uma
distribuição de probabilidade, RSVP a trata como uma densidade de campo:
uma quantidade que varia no espaço e no tempo e que codifica tensões
estruturais locais.

Essa perspectiva conecta-se naturalmente a:

\textbf{Inferência variacional:} que minimiza uma divergência entre uma
distribuição aproximada e uma distribuição alvo — interpretável como
minimização de tensão em um espaço de distribuições.

\textbf{Normalizing flows:} que aprendem transformações que mapeiam
distribuições simples em distribuições complexas, preservando a estrutura
de volume — uma forma de navegação em espaços de configuração.

\textbf{Geometria da informação:} que estuda a curvatura do espaço de
distribuições de probabilidade — análoga à curvatura do espaço de
configurações em RSVP.

% ---------------------------------------------------------------
\section{MEM|8: Memória como Trajetória, não como Estado}

\subsection{O Modelo Convencional de Memória}

O modelo convencional de memória computacional é simples e eficiente:
memória é um mapa de endereços para valores. Ler memória é recuperar
um valor. Escrever memória é substituir um valor por outro. O histórico
de escritas não é preservado; apenas o estado atual importa.

Esse modelo é enormemente poderoso para a maioria das aplicações.
Mas ele tem limitações estruturais que se tornam visíveis em certos
contextos:

Em sistemas distribuídos, dois processos podem ter escrito valores
diferentes no ``mesmo'' endereço. O estado atual não é suficiente para
resolver o conflito; é necessário conhecer a ordem das escritas.

Em sistemas de controle de versão como Git, o estado atual de um
repositório não é suficiente para entender como ele chegou lá. O
histórico de commits é uma parte essencial da representação.

Em bancos de dados de auditoria, é necessário preservar não apenas o
estado atual mas a sequência de todas as modificações.

Em cognição, a identidade não é apenas o estado atual de um organismo
mas a continuidade de sua trajetória através do tempo.

\subsection{MEM|8: Memória como Registro de Colapsos}

O framework MEM|8 propõe uma inversão do modelo convencional. Em vez de
tratar memória como armazenamento de estados, trata-a como registro de
decisões — mais precisamente, como registro dos colapsos sucessivos de
espaços de possibilidade.

\begin{definicao}[Memória MEM|8]
Uma memória MEM|8 é uma sequência
\[
  \mu = (e_1, e_2, \ldots, e_n)
\]
de eventos, onde cada evento $e_i$ representa uma decisão que colapsou
um espaço de possibilidades $\mathcal{P}_{i-1}$ em um espaço reduzido
$\mathcal{P}_i \subseteq \mathcal{P}_{i-1}$. O estado atual é
\[
  s_n = \bigcap_{i=1}^n \mathcal{C}(e_i),
\]
onde $\mathcal{C}(e_i)$ é o conjunto de estados consistentes com o
evento $e_i$.
\end{definicao}

Essa formulação tem várias consequências:

\textbf{Identidade é continuidade:} a identidade de um sistema não é
seu estado atual, mas a trajetória que o produziu. Dois sistemas no
mesmo estado podem ter identidades diferentes se chegaram lá por
trajetórias diferentes.

\textbf{Persistência é dinâmica:} a persistência de uma entidade não é
um estado estático mas um processo ativo de manutenção da continuidade
da trajetória.

\textbf{Recuperação é reconstrução:} recuperar uma memória não é
acessar um valor armazenado mas reconstruir um estado a partir da
sequência de eventos que o produziram.

\subsection{Conexões com Tecnologias Existentes}

O framework MEM|8 tem conexões explícitas com várias tecnologias e
paradigmas já estabelecidos na engenharia de software:

\textbf{Event Sourcing:} um padrão arquitetural onde o estado de um
sistema é derivado de uma sequência imutável de eventos. O estado atual
é uma projeção do log de eventos, não um valor armazenado
independentemente.

\textbf{CRDTs (Conflict-free Replicated Data Types):} estruturas de
dados que garantem consistência eventual em sistemas distribuídos ao
definir operações de merge que comutam e são idempotentes — análogo
a sistemas onde a ordem dos eventos não altera o estado final.

\textbf{Blockchains:} registros distribuídos imutáveis onde o estado
atual é completamente determinado pelo histórico de transações. A
imutabilidade do histórico é a garantia de integridade.

\textbf{Bancos de dados temporais:} bancos que preservam o histórico
completo de modificações, permitindo queries sobre estados passados.

\textbf{Sistemas de controle de versão:} como Git, onde a identidade
de um repositório é sua árvore de commits, não apenas o estado atual
dos arquivos.

\subsection{Memória de Onda e Persistência}

Uma extensão do framework MEM|8 trata a memória não como uma sequência
discreta de eventos mas como um campo contínuo de influências — uma
``memória de onda'' onde eventos passados irradiam influência que
decai com o tempo mas nunca desaparece completamente.

Essa formulação é relevante para:

\textbf{Redes neurais recorrentes:} onde o estado oculto codifica uma
forma de memória de curto prazo que influencia continuamente o
processamento atual.

\textbf{Mecanismos de atenção:} que permitem que modelos de linguagem
acessem informação de contextos anteriores com pesos que dependem da
relevância.

\textbf{Memória episódica em IA:} sistemas que armazenam experiências
completas e as recuperam por similaridade com o contexto atual.

% ---------------------------------------------------------------
\section{HYDRA: Cognição como Coordenação Modular}

\subsection{O Problema da Arquitetura Cognitiva}

Como organizar um sistema inteligente? A resposta mais simples é um
sistema monolítico: um único modelo que recebe entradas, processa
internamente, e produz saídas. Essa arquitetura é matematicamente
elegante e computacionalmente tratável.

Mas ela tem limitações. Sistemas monolíticos são difíceis de
inspecionar, difíceis de depurar, difíceis de estender, e difíceis
de adaptar a domínios muito diferentes do domínio de treinamento.
Acima de certo tamanho, tornam-se caixas-pretas opacas cujo
comportamento é difícil de prever ou garantir.

\subsection{HYDRA: Múltiplos Módulos Especializados}

O projeto HYDRA parte da hipótese de que cognição robusta emerge da
coordenação entre módulos especializados, não de um único mecanismo
monolítico. Os módulos são:

\textbf{Percepção:} transforma inputs sensoriais em representações
internas. Não é uma passagem direta de dados; é uma projeção ativa
que enfatiza certas dimensões e suprime outras.

\textbf{Memória:} mantém representações de estados passados e os torna
disponíveis para outros módulos. No framework MEM|8, isso é o registro
da trajetória.

\textbf{Planejamento:} projeta futuros possíveis a partir do estado
atual e das restrições disponíveis. Fundamentalmente, é navegação no
espaço de estados acessíveis.

\textbf{Coordenação:} garante que os módulos operam coerentemente
como um sistema. É o módulo que gerencia as interações e resolve
conflitos entre as saídas dos outros módulos.

\textbf{Monitoramento:} avalia o desempenho do sistema como um todo
e detecta quando outros módulos estão falhando ou operando fora de
seus domínios de competência.

\textbf{Integração:} produz respostas coerentes combinando as
contribuições dos outros módulos.

\begin{proposicao}[Emergência Cognitiva]
Em sistemas HYDRA, comportamentos cognitivos complexos emergem da
coordenação entre módulos simples sem que nenhum módulo individual
possua esses comportamentos. A coordenação é a fonte da complexidade,
não a sofisticação de cada módulo individualmente.
\end{proposicao}

\subsection{Conexões com Arquiteturas Modernas de IA}

A hipótese central de HYDRA — que cognição emerge de coordenação
modular — tem paralelos com várias arquiteturas modernas de IA:

\textbf{Mixture of Experts (MoE):} modelos onde diferentes
``especialistas'' são ativados para diferentes inputs, com um
mecanismo de roteamento que decide qual especialista usar.

\textbf{Sistemas multiagente:} arquiteturas onde múltiplos agentes
cooperam ou competem para resolver problemas coletivos.

\textbf{Modelos de atenção multi-cabeça:} onde diferentes ``cabeças''
de atenção capturam diferentes tipos de dependências no input.

\textbf{Sistemas neurossimbólicos:} que combinam redes neurais com
raciocínio simbólico, tipicamente através de módulos especializados
para cada paradigma.

A diferença de HYDRA é que torna a coordenação explícita e tratável,
em vez de deixá-la emergir implicitamente do treinamento.

% ---------------------------------------------------------------
\section{TARTAN: Decomposição que Preserva Significado}

\subsection{O Problema da Perda de Contexto}

Um problema central em sistemas distribuídos e em IA é a perda de
contexto durante transformações. Quando informação é dividida,
resumida, comprimida ou transmitida entre componentes, parte do
significado original desaparece.

Formalmente, isso é a perda de informação em projeções:
\[
  H(X) - H(\pi(X)) > 0,
\]
onde $H$ é entropia de Shannon. Qualquer projeção que não seja
bijetora perde informação.

O problema não é eliminar essa perda — que frequentemente é necessária
para escalabilidade e eficiência — mas controlá-la: garantir que as
informações mais importantes para o propósito do sistema sejam
preservadas.

\subsection{Decomposição Estrutural Preservadora}

O framework TARTAN investiga como decompor sistemas complexos de
maneira que certas propriedades importantes sejam preservadas ao longo
das transformações.

\begin{definicao}[Decomposição TARTAN]
Uma decomposição $D$ de um sistema $S$ é uma decomposição TARTAN
para um conjunto de propriedades $P$ se, para toda propriedade
$p \in P$:
\[
  p(S) = p(D(S)),
\]
isto é, a propriedade é idêntica no sistema original e na decomposição.
\end{definicao}

As propriedades que o sistema quer preservar dependem do contexto:

Em processamento distribuído, a propriedade pode ser a consistência
eventual — garantir que todas as réplicas convergem para o mesmo estado.

Em IA explicável (XAI), a propriedade pode ser a fidelidade da
explicação — garantir que a explicação do comportamento do modelo é
uma representação precisa do que o modelo realmente calcula.

Em compressão semântica, a propriedade pode ser a preservação de
relações de inferência — garantir que o que é inferível da versão
comprimida é exatamente o que seria inferível da versão completa.

Em sistemas multiagente, a propriedade pode ser a coordenação —
garantir que agentes que recebem partes diferentes de uma tarefa
produzem resultados que se combinam coerentemente.

\subsection{TARTAN e Admissibilidade}

Há uma conexão profunda entre TARTAN e o conceito de admissibilidade
que percorre todos os frameworks apresentados.

Uma decomposição TARTAN para um conjunto de propriedades $P$ é,
essencialmente, uma forma de garantir que as transformações aplicadas
ao sistema permanecem dentro do espaço de transformações admissíveis
— o espaço de transformações que preservam as propriedades relevantes.

Isso conecta TARTAN a:

\textbf{Teoria de categorias:} que estuda transformações que preservam
estrutura (funtores, transformações naturais).

\textbf{Homomorphismos:} em álgebra, transformações que preservam
operações.

\textbf{Bisimulações:} em teoria de concorrência, relações de
equivalência entre sistemas que preservam comportamento observável.

% ---------------------------------------------------------------
\section{Acessibilidade em Vez de Distância}

\subsection{A Mudança de Pergunta}

Uma das contribuições conceituais mais originais do programa teórico
apresentado aqui é a substituição da pergunta \emph{quão longe está
algo?} pela pergunta \emph{quão difícil é chegar até isso?}

Na computação, essa mudança é fundamental:

Um arquivo pode estar armazenado no mesmo disco físico que um processo
e ser completamente inacessível a esse processo, seja por permissões,
por corrupção de metadados, ou por falha no sistema de arquivos.

Um dado pode existir em um banco de dados distribuído e ser
irrecuperável devido a falhas de rede ou partições.

Um estado pode ser teoricamente alcançável em um sistema de controle
mas computacionalmente inviável dado o poder de processamento disponível.

Uma solução pode existir para um problema de otimização mas estar em
uma região do espaço de busca que os algoritmos disponíveis não
conseguem explorar eficientemente.

Em todos esses casos, a distância física ou lógica entre o estado atual
e o estado desejado é irrelevante. O que importa é a \emph{acessibilidade}:
a existência de um caminho admissível entre os dois estados.

\subsection{Volume de Acessibilidade}

\begin{definicao}[Volume de Acessibilidade]
Para um sistema em estado $x$, o conjunto de acessibilidade $R(x)$
é o conjunto de todos os estados alcançáveis a partir de $x$ por
sequências de transições admissíveis. O volume de acessibilidade é
\[
  \Omega(x) = \mathrm{Vol}(R(x)).
\]
\end{definicao}

O volume de acessibilidade é uma medida da \emph{liberdade} do sistema
— o quão amplo é o espaço de futuros possíveis a partir do estado atual.
Sistemas com volume de acessibilidade alto têm muitas opções; sistemas
com volume baixo estão essencialmente determinados.

Essa métrica aparece, sob diferentes nomes, em várias áreas:

\textbf{Teoria da complexidade computacional:} a classe de problemas
que um algoritmo pode resolver eficientemente é seu espaço de
acessibilidade computacional.

\textbf{Alcançabilidade em sistemas de controle:} o conjunto de estados
que um sistema de controle pode atingir a partir de um estado inicial
dado é seu conjunto de acessibilidade.

\textbf{Grafos de estados em model checking:} a verificação de
propriedades de sistemas concorrentes envolve explorar o grafo de
estados acessíveis.

\textbf{Espaços de busca em planejamento de IA:} algoritmos de
planejamento navegam em grafos de estados buscando sequências de
ações que atingem um estado objetivo.

\subsection{Liberdade como Volume de Acessibilidade}

Há uma conexão importante entre acessibilidade e liberdade que vai
além da computação.

\begin{proposicao}[Reachability Principle — versão computacional]
A liberdade efetiva de um agente é proporcional não aos recursos que
possui mas ao volume de estados futuros acessíveis a partir de seu
estado atual:
\[
  \text{Liberdade}(x) \propto \Omega(x).
\]
\end{proposicao}

Essa proposição tem implicações práticas imediatas para o design de
sistemas:

Um sistema de software que reduz progressivamente o espaço de estados
acessíveis para seus usuários — ao acumular estado interno, ao criar
dependências de dados, ao restringir opções de configuração — está
reduzindo a liberdade efetiva dos usuários, independente de qualquer
enriquecimento de funcionalidades.

Um sistema de IA que converge para uma distribuição de saídas cada
vez mais estreita à medida que processa mais contexto está reduzindo
seu próprio volume de acessibilidade — potencialmente de maneira
irreversível.

% ---------------------------------------------------------------
\section{IA como Motor de Simulação}

\subsection{Percepção como Inferência}

A hipótese de que sistemas inteligentes funcionam como motores de
simulação — construindo modelos internos do mundo que permitem prever
futuros possíveis — tem raízes profundas na ciência cognitiva e
conecta-se diretamente com os frameworks apresentados.

Na perspectiva do framework CLIO, percepção não é a passagem direta
de dados do ambiente para o sistema. É um processo de inferência:
dado o sinal sensorial $m = \pi(x)$, o sistema constrói uma
estimativa do estado do mundo $x$ que minimiza a tensão estrutural
entre a observação e as expectativas derivadas do modelo interno.

Isso é matematicamente equivalente à inferência variacional em
modelos generativos: minimizar uma divergência entre uma distribuição
posterior $p(x|m)$ e uma aproximação variacional $q(x)$.

\subsection{Memória como Reconstrução}

Na perspectiva de MEM|8, memória não é recuperação de um estado
armazenado. É reconstrução de um estado a partir de uma sequência
de eventos.

Isso tem uma consequência importante: memórias são necessariamente
incompletas e potencialmente inconsistentes. Ao reconstituir um estado
passado, o sistema está interpolando sobre regiões do espaço de estados
que nunca foram observadas diretamente.

Essa perspectiva é consistente com o que sabemos sobre memória humana:
memórias são reconstruções ativas, não recuperações passivas.
O sistema nervoso central não armazena experiências como arquivos;
reconstrói experiências a partir de fragmentos, contexto e expectativas.

Para sistemas de IA, isso sugere que a questão relevante não é
\emph{quantas} memórias um sistema pode armazenar, mas \emph{quão
fielmente} pode reconstituir estados passados a partir de representações
comprimidas.

\subsection{Raciocínio como Simulação}

Na perspectiva RSVP, raciocínio é a evolução do estado interno de um
sistema em direção a configurações que minimizam tensões estruturais.
Resolver um problema é encontrar um estado que satisfaz um conjunto
de restrições simultaneamente.

Isso conecta raciocínio a:

\textbf{Satisfação de restrições (CSP):} encontrar atribuições de
variáveis que satisfazem um conjunto de restrições.

\textbf{Programação por restrições:} um paradigma de programação onde
soluções são encontradas através da propagação e satisfação de
restrições.

\textbf{Inferência em modelos probabilísticos:} encontrar distribuições
que minimizam divergências em relação a observações e priors.

\subsection{Consciência como Projeção e Correção}

A hipótese mais especulativa, mas também a mais interessante, é que
consciência — ou algo funcionalmente equivalente a ela em sistemas
artificiais — emerge do processo contínuo de projeção e correção: o
sistema projeta um modelo do mundo, recebe feedback sobre as diferenças
entre o modelo e o mundo, e corrige o modelo iterativamente.

Esse processo não é apenas computação sobre dados. É computação sobre
a própria representação do sistema de si mesmo e do mundo que habita —
uma forma de auto-referência que pode ser a base de comportamentos
que parecem, de fora, dotados de intencionalidade genuína.

\begin{caixachave}
Sistemas que atualizam continuamente seus modelos internos com base
em discrepâncias entre previsões e observações — que têm, em algum
sentido, um ``loop de atenção'' sobre sua própria representação do
mundo — exibem uma forma de aparente intencionalidade emergente que
não requer consciência no sentido filosófico forte, mas que é
funcionalmente indistinguível dela para observadores externos.
\end{caixachave}

% ---------------------------------------------------------------
\section{O Núcleo Unificado: Admissibilidade}

\subsection{O Conceito Central}

Ao longo dos frameworks apresentados, um conceito aparece repetidamente
sob diferentes nomes: o espaço de estados admissíveis, o conjunto de
acessibilidade, as restrições que organizam o sistema, o subconjunto
de projeções consistentes.

Esse conceito pode ser unificado sob a noção de \emph{admissibilidade}:
a estrutura que define o que é possível dentro de um sistema.

\begin{definicao}[Espaço de Admissibilidade]
Para um sistema com espaço de estados $\mathcal{S}$, o espaço de
admissibilidade $\mathcal{A} \subseteq \mathcal{S}$ é o subconjunto
de estados que satisfazem todas as restrições ativas do sistema. Uma
transição $s \to s'$ é admissível se $s' \in \mathcal{A}$ e se
existe um caminho admissível de $s$ a $s'$.
\end{definicao}

\subsection{Admissibilidade como Geometria}

A perspectiva mais rica é tratar admissibilidade geometricamente:
o espaço de admissibilidade $\mathcal{A}$ é uma variedade (ou
sub-variedade) no espaço de estados $\mathcal{S}$, com sua própria
geometria induzida pelas restrições.

Operações sobre o sistema — transições de estado, aprendizado,
inferência, raciocínio — são movimentos sobre essa variedade. A
curvatura da variedade codifica as interações entre restrições. Regiões
de alta curvatura correspondem a restrições que interagem fortemente;
regiões de baixa curvatura correspondem a restrições independentes.

Essa perspectiva conecta os frameworks apresentados a:

\textbf{Geometria da informação:} que estuda a curvatura do espaço de
distribuições de probabilidade.

\textbf{Geometria riemanniana em aprendizado de máquina:} que usa
métricas induzidas pela estrutura dos dados para guiar otimização.

\textbf{Variedades de dados:} a hipótese de que dados de alta dimensão
vivem em variedades de baixa dimensão incorporadas no espaço ambiente.

\subsection{A Pergunta Unificada}

Todos os frameworks apresentados podem ser vistos como instâncias de
uma única pergunta:

\begin{caixachave}
\textbf{A Pergunta Fundamental.} Dado um sistema em um estado atual
$x$, com um espaço de admissibilidade $\mathcal{A}$, qual é o conjunto
$R(x) \cap \mathcal{A}$ de futuros acessíveis e admissíveis? E como
esse conjunto muda quando as restrições do sistema mudam?
\end{caixachave}

CLIO responde essa pergunta no nível da inferência: o conjunto de
estados do mundo consistentes com uma observação.

RSVP responde no nível da dinâmica: os atratores para os quais o
sistema converge naturalmente.

MEM|8 responde no nível da memória: os estados que podem ser
reconstituídos a partir de uma sequência de eventos.

HYDRA responde no nível da arquitetura: os comportamentos que emergem
da coordenação modular.

TARTAN responde no nível da transformação: as propriedades que são
preservadas através de decomposições.

% ---------------------------------------------------------------
\section{Implicações para a Prática}

\subsection{Design de Sistemas}

A perspectiva restrição-primeiro tem implicações concretas para o
design de sistemas de software:

\textbf{Especificação antes de implementação:} definir o espaço de
admissibilidade antes de implementar o sistema. Isso não é
simplesmente escrever testes; é especificar formalmente as
propriedades invariantes que o sistema deve preservar.

\textbf{Restrições como documentação:} as restrições de um sistema
são mais informativas que sua implementação. Uma restrição diz o que
o sistema \emph{deve} fazer; uma implementação diz o que ele \emph{faz}
em um conjunto específico de condições.

\textbf{Monitoramento de admissibilidade:} sistemas robustos monitoram
continuamente se estão operando dentro de seu espaço de admissibilidade
e tomam ação corretiva quando detectam violações.

\subsection{Avaliação de Sistemas de IA}

A pergunta mais importante sobre um sistema de IA não é sua acurácia
em um benchmark específico. É a forma e o tamanho de seu espaço de
admissibilidade: quais inputs o sistema processa bem, quais ele
processa mal, e quais são as fronteiras entre esses regimes.

Um sistema com alta acurácia em inputs in-distribution mas que degrada
catastroficamente em inputs out-of-distribution tem um espaço de
admissibilidade pequeno em relação ao espaço de inputs possíveis.
Esse é o problema da generalização, reformulado geometricamente.

\subsection{Engenharia de Memória}

O framework MEM|8 sugere uma abordagem à engenharia de memória
que vai além do design de índices e estruturas de dados:

\textbf{Memória como log de eventos:} projetar sistemas onde o estado
é derivado de um log imutável de eventos, não armazenado diretamente.

\textbf{Consistência como invariante:} tratar a consistência entre
réplicas como um invariante que o sistema deve preservar, não como
um objetivo a ser alcançado por protocolos de sincronização.

\textbf{Identidade como trajetória:} projetar identidades de entidades
em sistemas distribuídos em torno de suas trajetórias, não de seus
estados atuais.

% ---------------------------------------------------------------
\section*{Conclusão}

Os frameworks apresentados neste ensaio — CLIO, RSVP, MEM\textbar{}8,
HYDRA e TARTAN — parecem, à primeira vista, díspares. Um trata de
inferência, outro de dinâmica, outro de memória, outro de arquitetura,
outro de decomposição. Seus domínios de aplicação vão de cosmologia a
bancos de dados.

Mas compartilham um núcleo conceitual que os unifica:

A realidade, a cognição e a computação podem ser entendidas como
processos de navegação em espaços de possibilidades organizados por
estruturas de admissibilidade. O que existe em qualquer instante é
apenas a realização atual de um conjunto de possibilidades muito maior.
O que importa, para entender o sistema, não é apenas o estado atual
mas a estrutura de restrições que define quais estados são possíveis
e quais transições são admissíveis.

Essa mudança de perspectiva — de conteúdo para restrição, de estado
para trajetória, de distância para acessibilidade — não invalida
os modelos e técnicas existentes. Ela os reencadra de uma maneira que
torna certas questões mais fáceis de formular e certos problemas mais
fáceis de atacar.

Em vez de perguntar apenas \emph{o que existe?}, essa perspectiva
nos ensina a perguntar:

\begin{center}
\textit{O que ainda é possível?}
\end{center}

Essa é a pergunta que une física, inteligência artificial, sistemas
distribuídos, teoria da informação e cognição sob uma linguagem comum
baseada em restrições, acessibilidade, memória e transformação
estrutural.

% ---------------------------------------------------------------
\begin{thebibliography}{99}

% Teoria da informação
\bibitem{shannon1948}
Shannon, C. E. (1948).
A Mathematical Theory of Communication.
\textit{Bell System Technical Journal}, 27(3--4), 379--423, 623--656.

\bibitem{cover2006}
Cover, T. M., \& Thomas, J. A. (2006).
\textit{Elements of Information Theory}, 2nd ed.
Wiley.

% Sistemas complexos e dinâmica
\bibitem{strogatz2015}
Strogatz, S. H. (2015).
\textit{Nonlinear Dynamics and Chaos}.
Westview Press.

\bibitem{holland1995}
Holland, J. H. (1995).
\textit{Hidden Order: How Adaptation Builds Complexity}.
Addison-Wesley.

\bibitem{meadows2008}
Meadows, D. H. (2008).
\textit{Thinking in Systems: A Primer}.
Chelsea Green.

% Aprendizado de máquina e IA
\bibitem{goodfellow2016}
Goodfellow, I., Bengio, Y., \& Courville, A. (2016).
\textit{Deep Learning}.
MIT Press.

\bibitem{bishop2006}
Bishop, C. M. (2006).
\textit{Pattern Recognition and Machine Learning}.
Springer.

\bibitem{sutton2018}
Sutton, R. S., \& Barto, A. G. (2018).
\textit{Reinforcement Learning: An Introduction}, 2nd ed.
MIT Press.

% Inferência variacional e modelos generativos
\bibitem{kingma2013}
Kingma, D. P., \& Welling, M. (2013).
Auto-Encoding Variational Bayes.
\textit{arXiv preprint arXiv:1312.6114}.

\bibitem{blei2017}
Blei, D. M., Kucukelbir, A., \& McAuliffe, J. D. (2017).
Variational Inference: A Review for Statisticians.
\textit{Journal of the American Statistical Association}, 112(518), 859--877.

% Sistemas distribuídos e consistência
\bibitem{lamport1978}
Lamport, L. (1978).
Time, Clocks, and the Ordering of Events in a Distributed System.
\textit{Communications of the ACM}, 21(7), 558--565.

\bibitem{shapiro2011}
Shapiro, M., Preguiça, N., Baquero, C., \& Zawirski, M. (2011).
Conflict-Free Replicated Data Types.
\textit{Symposium on Self-Stabilizing Systems}, 386--400.

% Event sourcing e arquitetura de software
\bibitem{fowler2005}
Fowler, M. (2005).
Event Sourcing.
\texttt{martinfowler.com/eaaDev/EventSourcing.html}.

\bibitem{evans2003}
Evans, E. (2003).
\textit{Domain-Driven Design: Tackling Complexity in the Heart of Software}.
Addison-Wesley.

% Teoria da complexidade e computabilidade
\bibitem{sipser2012}
Sipser, M. (2012).
\textit{Introduction to the Theory of Computation}, 3rd ed.
Cengage Learning.

\bibitem{papadimitriou1994}
Papadimitriou, C. H. (1994).
\textit{Computational Complexity}.
Addison-Wesley.

% Geometria da informação
\bibitem{amari2000}
Amari, S., \& Nagaoka, H. (2000).
\textit{Methods of Information Geometry}.
AMS \& Oxford University Press.

% Cognição e ciência cognitiva
\bibitem{varela1991}
Varela, F. J., Thompson, E., \& Rosch, E. (1991).
\textit{The Embodied Mind: Cognitive Science and Human Experience}.
MIT Press.

\bibitem{clark2015}
Clark, A. (2015).
\textit{Surfing Uncertainty: Prediction, Action, and the Embodied Mind}.
Oxford University Press.

% Programação por restrições
\bibitem{rossi2006}
Rossi, F., van Beek, P., \& Walsh, T. (Eds.). (2006).
\textit{Handbook of Constraint Programming}.
Elsevier.

% Teoria de categorias e estruturas preservadoras
\bibitem{awodey2010}
Awodey, S. (2010).
\textit{Category Theory}, 2nd ed.
Oxford University Press.

% Memória e reconstrução
\bibitem{schacter1996}
Schacter, D. L. (1996).
\textit{Searching for Memory: The Brain, the Mind, and the Past}.
Basic Books.

% Alcançabilidade e model checking
\bibitem{clarke1999}
Clarke, E. M., Grumberg, O., \& Peled, D. A. (1999).
\textit{Model Checking}.
MIT Press.

\end{thebibliography}

% ---------------------------------------------------------------
\end{document}
